\documentclass[11pt]{article}
\usepackage[margin=1in]{geometry}
\usepackage{amsmath, amssymb, amsthm}
\usepackage{hyperref}
\usepackage{booktabs}

\title{Sampling-noise floor of the Jensen--Shannon distance \\
       between two empirical histograms}
\author{}
\date{}

\newtheorem{proposition}{Proposition}

\begin{document}
\maketitle

\section{Setup}

Let $P$ be a fixed distribution on $K$ bins with PMF $p_1,\dots,p_K$
($\sum_k p_k = 1$). Draw two independent samples $D_1,D_2$ of size
$N$ each from $P$, and let
\[
  \hat p_k = \frac{1}{N}\sum_{x\in D_1}\mathbf{1}\{x\in \text{bin }k\},
  \qquad
  \hat q_k = \frac{1}{N}\sum_{x\in D_2}\mathbf{1}\{x\in \text{bin }k\}
\]
be the empirical PMFs. Define the midpoint
$\hat m_k = (\hat p_k + \hat q_k)/2$ and the
Jensen--Shannon divergence (in bits)
\[
  \mathrm{JS}(\hat p,\hat q)
  \;=\;
  \tfrac12\,\mathrm{KL}(\hat p \,\|\, \hat m)
  +\tfrac12\,\mathrm{KL}(\hat q \,\|\, \hat m),
\]
and the JS \emph{distance} $d \;=\; \sqrt{\mathrm{JS}(\hat p,\hat q)}\in[0,1]$.

Our question: under the null hypothesis that $D_1$ and $D_2$ come from the
same distribution, what is $\mathbb{E}[d]$?

\section{Result}

\begin{proposition}\label{prop:main}
For $N$ large and the two samples drawn iid from a common distribution
with $K$ bins,
\[
  \boxed{\;\mathbb{E}[d]
  \;\approx\;
  \sqrt{\dfrac{K-\tfrac32}{4\,N\,\ln 2}}\;}.
\]
\end{proposition}

For $N=1000$ and $K=30$, this gives
\(
\sqrt{28.5 / (4000\ln 2)} \approx 0.101,
\)
which matches the empirical floor we observe in our experiments
(values in the range $0.10$--$0.12$ across queries).

\section{Proof}

\paragraph{Step 1: quadratic expansion of JS divergence.}
Write $\delta_k = \hat p_k - \hat q_k$, so that
$\hat p_k - \hat m_k = \delta_k/2$ and
$\hat q_k - \hat m_k = -\delta_k/2$. A second-order Taylor expansion of
the KL divergence around its second argument gives
\[
  \mathrm{KL}(r \,\|\, s) \;=\; \frac{1}{\ln 2}\sum_k
  \left[(r_k-s_k) + \tfrac{(r_k-s_k)^2}{2s_k} + O\!\big((r_k-s_k)^3\big)\right],
\]
where the linear term drops out because $\sum_k r_k=\sum_k s_k=1$. Hence
\begin{align*}
  \mathrm{KL}(\hat p \,\|\, \hat m) +
  \mathrm{KL}(\hat q \,\|\, \hat m)
  &\approx
  \frac{1}{2\ln 2}
  \sum_k \frac{(\delta_k/2)^2 + (\delta_k/2)^2}{\hat m_k}
  \\&=\;
  \frac{1}{4\ln 2}\sum_k \frac{\delta_k^2}{\hat m_k},
\end{align*}
and therefore
\begin{equation}\label{eq:quadratic}
  \mathrm{JS}(\hat p,\hat q)
  \;\approx\;
  \frac{1}{8\ln 2}\sum_k \frac{\delta_k^2}{\hat m_k}.
\end{equation}

\paragraph{Step 2: expectation.}
Under the null, $\mathbb{E}[\delta_k]=0$ and $\hat m_k \to p_k$ in
probability. Moreover, $\hat p_k$ and $\hat q_k$ are independent, each
$\mathrm{Binomial}(N,p_k)/N$, so
\[
  \mathbb{E}[\delta_k^2] \;=\; \mathrm{Var}(\hat p_k) + \mathrm{Var}(\hat q_k)
  \;=\; \frac{2\,p_k(1-p_k)}{N}.
\]
Plugging into~\eqref{eq:quadratic},
\begin{equation}\label{eq:EJS}
  \mathbb{E}[\mathrm{JS}]
  \;\approx\;
  \frac{1}{8\ln 2}\sum_k \frac{2\,p_k(1-p_k)/N}{p_k}
  \;=\;
  \frac{1}{4N\ln 2}\sum_k (1-p_k)
  \;=\;
  \frac{K-1}{4N\ln 2}.
\end{equation}
This expectation depends only on $K$ and $N$ --- not on the shape of $P$.

\paragraph{Step 3: from divergence to distance.}
The normalized statistic
\[
  T \;=\; \frac{N}{2}\sum_k \frac{\delta_k^2}{p_k}
\]
is asymptotically $\chi^2_{K-1}$ distributed: the differences
$\delta_k/\sqrt{2p_k/N}$ are approximately standard normal and satisfy
the single linear constraint $\sum_k \delta_k = 0$. Combining with
\eqref{eq:quadratic} and $\hat m_k \to p_k$,
\[
  \mathrm{JS}\;\approx\; \frac{T}{4N\ln 2},
\]
so
\[
  d \;=\; \sqrt{\mathrm{JS}} \;\approx\;
  \frac{\sqrt{T}}{\sqrt{4N\ln 2}}
  \;\sim\;
  \frac{\sqrt{\chi^2_{K-1}}}{\sqrt{4N\ln 2}}.
\]
For moderate--large $K$, the chi distribution concentrates with
\(\mathbb{E}\!\left[\sqrt{\chi^2_\nu}\right] \approx \sqrt{\nu-\tfrac12}\),
hence
\[
  \mathbb{E}[d] \;\approx\; \sqrt{\frac{K-\tfrac32}{4N\ln 2}}.
\]
\qed

\section{Discussion}

\paragraph{Scaling.}
The floor decays as $1/\sqrt{N}$ and grows as $\sqrt{K}$. Doubling $N$
shrinks the typical JS distance by $\sqrt{2}$; doubling the bin count
inflates it by roughly $\sqrt{2}$.

\paragraph{Numerical values.}
\begin{center}
\begin{tabular}{rrl}
\toprule
$N$ & $K$ & $\mathbb{E}[d]$ \\
\midrule
$1000$  & $20$  & $0.083$ \\
$1000$  & $30$  & $0.101$ \\
$1000$  & $100$ & $0.188$ \\
$10000$ & $30$  & $0.032$ \\
\bottomrule
\end{tabular}
\end{center}

\paragraph{Practical implication.}
A JS distance near the value in Proposition~\ref{prop:main} is
consistent with the two samples coming from the same distribution;
nothing can be concluded from it.  A useful per-query test is the
two-sample Kolmogorov--Smirnov statistic on the raw samples (no
binning), or a permutation test on Wasserstein distance.

\end{document}
